\title{DeepSeekMath: Pushing the Limits of Mathematical Reasoning in Open Language Models}

\begin{abstract}
The intricacies and formal structure inherent in mathematical reasoning present notable difficulties for language models. This study presents DeepSeekMath~7B, an extension of DeepSeek-Coder-Base-v1.5 7B, further trained using 120B mathematics-focused tokens extracted from Common Crawl, which also includes linguistic and programming material. Without the aid of auxiliary toolkits or ensembling strategies, DeepSeekMath~7B attains a remarkable 51.7\% accuracy on the advanced-level MATH benchmark, nearing the results of state-of-the-art systems such as Gemini-Ultra and GPT-4. When leveraging self-consistency sampling with 64 outputs, performance on MATH rises to 60.9\%. The advanced mathematical reasoning observed in DeepSeekMath~originates primarily from two innovations: firstly, a robust data selection framework that maximally exploits the wealth of publicly available web resources; and secondly, the development of Group Relative Policy Optimization (GRPO)—an adaptation of Proximal Policy Optimization (PPO)—which not only bolsters reasoning capabilities in mathematics but also achieves improved efficiency in PPO's memory consumption.
\end{abstract}

\section{Introduction}

Large language models (LLM) have fundamentally transformed methods of mathematical reasoning within artificial intelligence, leading to marked progress on both quantitative reasoning tasks \citep{MATH} and geometry-focused benchmarks \citep{trinh2024solving}. Additionally, these models have shown considerable utility in supporting human problem-solvers tackling intricate mathematical challenges \citep{tao}. Nevertheless, state-of-the-art systems like GPT-4 \citep{gpt4} and Gemini-Ultra \citep{gemini} remain closed to the public, while existing open-source alternatives still lag significantly in terms of capabilities and overall performance.

In this work, we present DeepSeekMath, a specialized language model that markedly advances the mathematical reasoning capabilities of open-source models and achieves near GPT-4-level results on several academic benchmarks. To facilitate this, we construct the DeepSeekMath Corpus, an extensive and high-quality dataset for pre-training, consisting of 120B math-focused tokens. Our dataset is sourced from the Common Crawl (CC) through the application of a fastText-based classifier \citep{joulin2016fasttext}. Initially, the classifier is trained with positive samples drawn from OpenWebMath \citep{openwebmath} and a varied array of unrelated web pages as negative samples. After this first stage, we utilize the trained classifier to extract additional positive math content from the CC, which undergoes further verification and refinement by human annotators. This curated set is then employed to further train and enhance the classifier's accuracy. Evaluation demonstrates that our large-scale dataset maintains a high degree of quality, as reflected by DeepSeekMath-Base 7B achieving 64.2\% on GSM8K \citep{gsm8k} and 36.2\% on the MATH competition dataset \citep{MATH}, surpassing Minerva 540B \citep{minerva}. Furthermore, the multilingual composition of the DeepSeekMath Corpus results in observable gains on Chinese mathematics benchmarks \citep{wei2023cmath,agieval}. We view our methodology for mathematical data preparation as a valuable foundation for future research, anticipating considerable potential for further advancements.

DeepSeekMath-Base is built upon the DeepSeek-Coder-Base-v1.5 7B model \citep{deepseek-coder}, based on our observation that initializing from a code-focused model yields superior results relative to beginning with a standard LLM. In addition, our results reveal that mathematical training augments performance not only on mathematical tasks but also on broader reasoning evaluations such as MMLU \citep{mmlu} and BBH \citep{bbh} benchmarks. This suggests that such training strengthens both mathematical proficiency and general reasoning ability.

Following pre-training, we conduct mathematical instruction tuning on DeepSeekMath-Base using datasets involving chain-of-thought \citep{cot}, program-of-thought \citep{pot,pal}, and tool-integrated reasoning \citep{tora} methodologies. The finalized DeepSeekMath-Instruct 7B model surpasses other models of similar 7B scale and demonstrates performance on par with open-source instruction-tuned models at the 70B parameter level.

In addition, we propose Group Relative Policy Optimization (GRPO), a modified reinforcement learning (RL) algorithm derived from Proximal Policy Optimization (PPO) \citep{schulman2017proximal}. Unlike traditional approaches, GRPO dispenses with the critic model, instead leveraging group scores to estimate the baseline, which considerably lowers the computational demands of training. Utilizing only a portion of English instruction tuning data, GRPO demonstrates marked gains over the robust DeepSeekMath-Instruct model in both in-domain (GSM8K: 82.9\% $\rightarrow$ 88.2\%, MATH: 46.8\% $\rightarrow$ 51.7\%) and out-of-domain mathematical benchmarks (such as CMATH: 84.6\% $\rightarrow$ 88.8\%) during reinforcement learning. We further present a cohesive framework for interpreting diverse methods, encompassing Rejection Sampling Fine-Tuning (RFT) \citep{yuan2023scaling}, Direct Preference Optimization (DPO) \citep{dpo}, PPO, and GRPO. Within this unified framework, we identify that these approaches fundamentally represent either direct or streamlined forms of RL. Our study includes comprehensive experimental comparisons—such as online versus offline training, supervision at the level of outcomes versus processes, and single-step versus iterative RL—to thoroughly probe the critical features of this paradigm. Finally, we elucidate the mechanisms by which our RL methodology enhances the capabilities of instruction-tuned models and offer a synthesis of promising avenues for advancing RL based on our unified framework.

\subsection{Contributions}
We present scalable approaches to mathematical pre-training and conduct a comprehensive investigation into reinforcement learning methodologies.

\textbf{Math Pre-Training at Scale}

\begin{itemize}[topsep=0pt]
    \item In this study, we demonstrate that the openly available Common Crawl dataset holds significant resources pertinent to mathematics. Through a carefully crafted pipeline for data selection, we assembled the DeepSeekMath Corpus—a curated, high-caliber dataset comprised of 120B tokens extracted from web pages screened for mathematical relevance. This collection substantially surpasses previous datasets, being nearly sevenfold larger than the set of math web pages utilized by Minerva \citep{minerva} and approximately nine times larger than the newly published OpenWebMath \citep{openwebmath}.

\item DeepSeekMath-Base 7B, our pre-trained foundational model, matches the performance of Minerva 540B \citep{minerva}. This suggests that sheer parameter count is not the sole determinant of mathematical reasoning proficiency; robust results can also be attained by smaller models when trained on high-quality datasets.

\item We present the outcomes of our experiments involving math training. Conducting code training before introducing math training enhances models' proficiency in addressing mathematical tasks, regardless of whether tools are utilized. This provides partial insight into the enduring question: \textit{does code training improve reasoning abilities?} Our results suggest that it does, specifically with respect to mathematical reasoning.

\item While it is typical to utilize arXiv papers for training—particularly in mathematics-focused research—this approach does not yield significant gains across any of the mathematical benchmarks evaluated in this study.
\end{itemize}

\textbf{Exploration and Analysis of Reinforcement Learning}

\begin{itemize}[topsep=0pt]
    \item We propose Group Relative Policy Optimization (GRPO), a reinforcement learning technique that is both resource-efficient and robust. Unlike Proximal Policy Optimization (PPO), GRPO eliminates the need for a critic model by deriving the baseline directly from group scores, thereby lowering computational demands during training.
    \item Our results show that applying GRPO to instruction-tuning data substantially improves the performance of our instruction-tuned model, DeepSeekMath-Instruct. Additionally, we observe notable gains in generalization to out-of-domain tasks throughout reinforcement learning.
    \item We establish a cohesive framework to interpret various methodologies including RFT, DPO, PPO, and GRPO. Through a comprehensive set of experiments—comparing, for example, online versus offline training, outcome versus process supervision, and single-turn versus iterative reinforcement learning—we systematically analyze the key factors underpinning this framework.
    \item Leveraging this unified framework, we investigate the underlying mechanisms contributing to the success of reinforcement learning, and identify several promising avenues for enhancing reinforcement learning in large language models (LLMs).
\end{itemize}

\subsection{Summary of Evaluations and Metrics}

\begin{itemize}[topsep=0pt]
    \item \textbf{English and Chinese Mathematical Reasoning}: Our study systematically evaluates model performance on both English and Chinese mathematical reasoning datasets, encompassing questions ranging from elementary to university complexity. English evaluation sets feature GSM8K \citep{gsm8k}, MATH \citep{MATH}, SAT \citep{llemma}, OCW Courses \citep{minerva}, and MMLU-STEM \citep{mmlu}, while Chinese datasets comprise MGSM-zh \citep{mgsm}, CMATH \citep{wei2023cmath}, Gaokao-MathCloze \citep{agieval}, and Gaokao-MathQA \citep{agieval}. The evaluation focuses on models' proficiency in composing coherent, self-explanatory textual solutions independently of external tools, in addition to their capacity to address problems by leveraging Python code.

On English evaluation tasks, DeepSeekMath-Base achieves performance on par with the proprietary Minerva 540B model \citep{minerva} and consistently outperforms all available open-source base models, such as Mistral 7B \citep{mistral} and Llemma-34B \citep{llemma}, regardless of their involvement in mathematical pre-training, often by a substantial margin. Importantly, DeepSeekMath-Base exhibits clear superiority on Chinese assessments, which can be attributed to our approach of not limiting the pre-training corpus to English mathematics texts, as done in previous research \citep{minerva,llemma}, and by incorporating high-quality mathematical data in other languages. Building on this foundation with mathematical instruction tuning and reinforcement learning, our DeepSeekMath-Instruct and DeepSeekMath-RL models achieve notable results, surpassing 50\% accuracy on the challenging competition-grade MATH dataset—setting a new milestone among open-source solutions.

\item \textbf{Formal Mathematics}: We assess DeepSeekMath-Base on the informal-to-formal theorem proving challenge as proposed in \citep{dsp_proof}, utilizing the miniF2F benchmark \citep{minif2f} and employing Isabelle \citep{isabelle} as the proof assistant. The model achieves impressive few-shot results in autoformalization tasks.
\item \textbf{Natural Language Understanding, Reasoning, and Code}: To thoroughly gauge the models' proficiency in language comprehension, reasoning, and coding, we conduct evaluations of DeepSeekMath-Base with the Massive Multitask Language Understanding (MMLU) benchmark \citep{mmlu}, encompassing 57 multiple-choice problems across a wide variety of disciplines. Additionally, we use BIG-Bench Hard (BBH) \citep{bbh}, containing 23 complex tasks that primarily necessitate multi-step reasoning, along with HumanEval \citep{codex} and MBPP \citep{mbpp}, both standard benchmarks for code evaluation. Pre-training on mathematical data enhances both reasoning and language understanding capabilities.

\section{Math Pre-Training}

\subsection{Data Collection and Decontamination} This section describes the methodology used to build the DeepSeekMath Corpus utilizing Common Crawl data. As illustrated in Figure \ref{fig:data_collect}, we introduce an iterative workflow for efficiently compiling a substantial mathematical dataset from Common Crawl, beginning with a seed set that consists of a limited yet high-quality assembly of mathematics-focused data. Importantly, this framework can be readily extended to domains beyond mathematics, including areas like coding.

Our initial step involves selecting OpenWebMath \citep{openwebmath}, a well-curated collection of mathematical content sourced from the web, as our primary seed corpus. With this dataset, we train a fastText classifier \citep{joulin2016fasttext} to efficiently retrieve additional mathematical pages similar in nature to those within OpenWebMath. For this purpose, we sample 500,000 documents from the seed as positive instances and pair them with 500,000 negative examples, which are randomly selected web pages from Common Crawl. The model is trained with an open-source fastText implementation\footnote{\url{https://fasttext.cc}}, setting the embedding size to 256, the learning rate at 0.1, the maximum word n-gram size at 3, a word occurrence threshold of 3, and using 3 epochs. To effectively downsize the Common Crawl dataset, we apply deduplication based on URLs and implement near-duplicate filtering, thus narrowing it to approximately 40B HTML web documents. Next, we apply the trained fastText model to identify mathematical content from this reduced Common Crawl set. For quality control, pages are ordered by their fastText-predicted scores, retaining only the highest-scoring mathematical documents. The retained data volume is further examined via pre-training trials conducted on the top 40B, 80B, 120B, and 160B tokens. For this initial cycle, we opt to preserve the leading 40B tokens.

Following the initial round of data gathering, a substantial number of mathematical web pages remain uncaptured, primarily because the fastText model’s initial set of positive samples is not sufficiently varied. To enhance the model’s effectiveness, we expand the seed corpus by identifying and incorporating additional mathematical web sources. Our procedure begins by segmenting the entire Common Crawl into non-overlapping domains, where each domain corresponds to web pages under a single base URL. For each domain, we compute the proportion of web pages retrieved during the first iteration. Domains exhibiting more than 10\% collection coverage are considered to be mathematics-related (such as \url{mathoverflow.net}). Next, we manually review and label URLs within these domains that are associated with mathematical materials (for instance, \url{mathoverflow.net/questions}). Web pages associated with these URLs that were previously omitted are subsequently added to the seed corpus. This strategy allows us to increase the quantity of positive training examples and train a more robust fastText model, improving our ability to identify mathematical web content in future rounds. After performing four rounds of data collection, we ultimately accumulate 35.5M mathematical web pages, corresponding to 120B tokens. During the fourth round, we observe that approximately 98\% of the pages had already been acquired in the third round, prompting us to discontinue further data collection.

In order to prevent benchmark contamination, we adopt the filtering strategy proposed by \cite{deepseek-coder}. This involves excluding web pages that feature questions or answers from established English mathematical benchmarks like GSM8K \citep{gsm8k} and MATH \citep{MATH}, as well as Chinese benchmarks such as CMATH \citep{wei2023cmath} and AGIEval \citep{agieval}. Specifically, our approach removes any segment of text within the math training corpus if it contains a contiguous sequence of 10 words that precisely corresponds to any sub-string found in the evaluation benchmarks. For benchmark samples with fewer than 10 but at least 3 grams, we apply exact string matching to filter out the affected web pages from our dataset.

\subsection{Validating the Quality of the DeepSeekMath~Corpus}

We run pre-training experiments to investigate how the DeepSeekMath~Corpus is compared with the recently released math-training corpora:

\begin{itemize}[topsep=0pt]
    \item \textbf{MathPile} \citep{mathpile}: This dataset is composed of 8.9B tokens sourced from various platforms, such as textbooks, Wikipedia, ProofWiki, CommonCrawl, StackExchange, and arXiv; notably, arXiv accounts for more than 85\% of the total content;
    \item \textbf{OpenWebMath} \citep{openwebmath}: This corpus, comprising 13.6B tokens, is curated by extracting mathematical data from CommonCrawl using specialized filtering techniques;
    \item \textbf{Proof-Pile-2} \citep{llemma}: This collection integrates multiple resources, including OpenWebMath, AlgebraicStack (10.3B tokens of mathematical code), and arXiv literature (28.0B tokens). For experiments involving Proof-Pile-2, we adhere to the methodology of \cite{llemma}, employing an arXiv:Web:Code distribution of 2:4:1.
\end{itemize}

\subsubsection{Training Setting}
\label{sec:quality-policy}
We fine-tune a general-purpose language model, equipped with 1.3B parameters and architecturally aligned with the DeepSeek LLMs \citep{deepseek-llm}, henceforth referred to as DeepSeek-LLM 1.3B, for mathematical tasks. Each version of the model is optimized on an individual mathematical dataset for a total of 150B tokens. For all experimental runs, we utilize the streamlined and resource-efficient HAI-LLM \citep{haillm} training infrastructure. Mirroring the protocol employed for DeepSeek LLMs, training is performed with the AdamW optimizer \citep{adamW}, setting $\beta_1=0.9$, $\beta_2=0.95$, and $\mathrm{weight\_decay}=0.1$. A multi-phase learning rate strategy is adopted: the learning rate is gradually increased and attains its peak following 2,000 warmup iterations, then is reduced to 31.6\% of this maximum at 80\% training completion, and further diminished to 10.0\% of the peak value at the 90\% mark. The highest value for the learning rate is fixed at 5.3e-4. We select a global batch size of 4M tokens and employ a context window of 4K tokens throughout training.

\subsubsection{Evaluation Results}

\textbf{The DeepSeekMath~Corpus is of high quality, covers multilingual mathematical content, and is the largest in size.}

\begin{itemize}[topsep=0pt]
    \item \textbf{High-quality}:
    To assess downstream effectiveness, we test on 8 mathematical benchmark datasets utilizing few-shot chain-of-thought prompting \cite{cot}. Results in Table \ref{tab:corpora_comparison} reveal that models trained with the DeepSeekMath~Corpus substantially outperform competitors. Moreover, Figure \ref{fig:corpora_comparisons} illustrates that after processing 50B tokens (equivalent to one epoch of Proof-Pile-2), DeepSeekMath-trained models yield superior outcomes compared to those trained with Proof-Pile-2, suggesting that the DeepSeekMath~Corpus possesses a higher overall data quality.
    \item \textbf{Multilingual}:
    The DeepSeekMath~Corpus contains content in several languages, with English and Chinese making up the largest shares. According to Table \ref{tab:corpora_comparison}, incorporating the DeepSeekMath~Corpus into model training boosts mathematical reasoning capabilities in both English and Chinese. By comparison, other mathematical datasets, which largely focus on English-language material, tend to offer limited benefits for Chinese-language tasks and may in some cases negatively affect performance.
    \item \textbf{Large-scale}:
    In terms of scale, DeepSeekMath~Corpus significantly surpasses the size of previous mathematical corpora. As shown in Figure \ref{fig:corpora_comparisons}, DeepSeek-LLM 1.3B, after training on the DeepSeekMath~Corpus, demonstrates a more pronounced and sustained improvement. Meanwhile, smaller baseline datasets, which are reused throughout training, tend to rapidly approach a performance ceiling.
\end{itemize}

\subsection{Training and Evaluating DeepSeekMath-Base 7B}

This section presents DeepSeekMath-Base 7B, a foundational model exhibiting robust reasoning capabilities with a particular emphasis on mathematical problem solving. The model begins from DeepSeek-Coder-Base-v1.5 7B \citep{deepseek-coder} and undergoes pretraining over 500B tokens. The composition of the training data is as follows: DeepSeekMath Corpus constitutes 56\%, AlgebraicStack provides 4\%, arXiv accounts for 10\%, Github code makes up 20\%, while the final 10\% consists of natural language data sourced from Common Crawl in both English and Chinese. We largely adhere to the training regimen outlined in Section \ref{sec:quality-policy}; however, the maximum learning rate is set to 4.2e-4, and the batch size used is 10M tokens.

This study undertakes an extensive evaluation of DeepSeekMath-Base 7B's mathematical proficiency, examining its effectiveness in independently generating complete mathematical solutions, employing tools to address mathematical problems, and executing formal theorem proving. In addition to its mathematical assessment, we also offer a broader characterization of the base model, assessing its capabilities in natural language comprehension, logical reasoning, and programming.

\paragraph{Mathematical Problem Solving with Step-by-Step Reasoning}

We assess DeepSeekMath-Base's ability to solve mathematical tasks via few-shot chain-of-thought prompting \citep{cot} on eight different benchmarks presented in both English and Chinese. These evaluation sets span a range of quantitative reasoning tasks, such as GSM8K \citep{gsm8k}, MATH \citep{MATH}, and CMATH \citep{wei2023cmath}, as well as multiple-choice assessments like MMLU-STEM \citep{mmlu} and Gaokao-MathQA \citep{agieval}. Collectively, the selected benchmarks represent a broad spectrum of mathematical disciplines, from elementary concepts to college-level challenges.

According to Table \ref{tab:base_math_cot}, DeepSeekMath-Base 7B achieves the top scores on all eight benchmarks among open-source base models, including both the well-known Mistral 7B \citep{mistral} and the newer Llemma 34B \citep{llemma}, the latter of which received specialized mathematical training on Proof-Pile-2 \citep{llemma}. In particular, DeepSeekMath-Base demonstrates a greater than 10\% absolute improvement over previous open-source base models on the challenging MATH competition dataset. Furthermore, it surpasses the performance of the considerably larger closed-source Minerva 540B \citep{minerva}—which, at 77 times the size, builds on PaLM \citep{palm} with additional mathematical pretraining.

\paragraph{Mathematical Problem Solving with Tool Use} We assess the capability of models to perform mathematical reasoning augmented by external tools on the GSM8K and MATH benchmarks, employing a few-shot program-of-thought prompting framework as proposed in \citep{pot,pal}. In this setting, models address each mathematical query by generating a Python script, leveraging libraries like \textit{math} and \textit{sympy} for handling complex calculations. The computed output of the generated program serves as the final answer. According to the results in Table \ref{tab:base_math_pot_proof}, DeepSeekMath-Base 7B achieves superior performance over the previously leading Llemma 34B model.

\paragraph{Formal Mathematics}

Automating formal proofs offers significant advantages in bolstering the correctness and dependability of mathematical arguments, while also improving productivity—a trend that has garnered growing interest in recent years. In our study, we assess the DeepSeekMath-Base 7B model on the informal-to-formal proving task introduced in \citep{dsp_proof}, which involves generating a formal proof from an informal problem description, a matching formal statement, and an accompanying informal argument. Our evaluation uses the miniF2F dataset \citep{minif2f}, which serves as a benchmark for formalized Olympiad-grade mathematical problems. For each instance, we utilize few-shot prompting to guide the generation of formal proofs in Isabelle. Consistent with the methodology in \cite{dsp_proof}, we employ the model to produce an initial outline or sketch of the proof, then invoke the widely-used automated prover Sledgehammer \citep{sledgehammer} to complete any missing intermediate steps. Table \ref{tab:base_math_pot_proof} displays the outcomes, where DeepSeekMath-Base 7B achieves notable success in the task of proof autoformalization.

\paragraph{Natural Language Understanding, Reasoning, and Code}

We assess natural language understanding using MMLU \citep{mmlu}, reasoning abilities through BBH \citep{bbh}, and programming competence with the HumanEval \citep{codex} and MBPP \citep{mbpp} benchmarks. According to Table \ref{tab:understanding_reasoning_code}, DeepSeekMath-Base 7B achieves marked performance improvements on both MMLU and BBH compared to its predecessor, DeepSeek-Coder-Base-v1.5 \citep{deepseek-coder}, underscoring the beneficial effects of mathematical pretraining for tasks involving language comprehension and logical reasoning. Furthermore, the incorporation of code tokens in the continual training process enables DeepSeekMath-Base 7B to retain the coding performance demonstrated by DeepSeek-Coder-Base-v1.5 on HumanEval and MBPP. In sum, DeepSeekMath-Base 7B substantially surpasses the generalist model Mistral 7B \citep{mistral} across all three evaluated reasoning and coding benchmarks.

\section{Supervised Fine-Tuning}

\subsection{SFT Data Curation}

We develop an instruction-tuning dataset for mathematics that encompasses English and Chinese questions across a range of topics and difficulty levels. Each problem is accompanied by its respective solution formatted in chain-of-thought (CoT) \citep{cot}, program-of-thought (PoT) \citep{pot,pal}, and tool-integrated reasoning approaches \citep{tora}. In total, the training set consists of 776K examples.

\begin{itemize}[topsep=0pt]
    \item \textbf{English mathematical datasets}: We provide tool-assisted solution annotations for both GSM8K and MATH datasets, and further incorporate a selection from MathInstruct \citep{MathInstruct}, together with the Lila-OOD training data \citep{lila}, in which problems are resolved using CoT or PoT methodologies. The assembled English dataset spans a wide range of mathematical domains, including algebra, probability, number theory, calculus, and geometry.
    \item \textbf{Chinese mathematical datasets}: Our collection features K-12 mathematical problems in Chinese, distributed over 76 distinct sub-categories, such as linear equations. Each problem includes solutions with both CoT and tool-based reasoning annotations.
\end{itemize}

\subsection{Training and Evaluating DeepSeekMath-Instruct 7B}

This section presents DeepSeekMath-Instruct 7B, a model that is subjected to mathematical instruction tuning using DeepSeekMath-Base as its foundation. Training data samples are combined in a random sequence until they reach a total context window of 4K tokens. The model is trained over 500 iterations, employing a batch size of 256 and maintaining a steady learning rate of 5e-5 throughout the process.

We evaluate models' mathematical performance both without and with tool use, on 4 quantitative reasoning benchmarks in English and Chinese. We benchmark our model against the leading models of the time:

\begin{itemize}[topsep=0pt]
    \item \textbf{Closed-source models} comprise the following: (1) the GPT series, with GPT-4 \citep{gpt4} and the GPT-4 Code Interpreter~\footnote{\url{https://openai.com/blog/chatgpt-plugins\#\#code-interpreter}} standing out as the most advanced, (2) Gemini Ultra and Pro \citep{gemini}, (3) Inflection-2 \citep{inflection-2}, (4) Grok-1~\footnote{\url{https://x.ai/model-card}}, and additionally, some recent offerings from Chinese companies such as (5) Baichuan-3~\footnote{\url{https://www.baichuan-ai.com}}, and (6) the newest version of GLM, GLM-4~\footnote{\url{https://open.bigmodel.cn/dev/api\#glm-4}}

    from the GLM family \citep{glm}.

These models are intended for broad applications, and the majority have been subjected to various alignment techniques.

\item \textbf{Open-source models} comprise: general-purpose models such as (1) DeepSeek-LLM-Chat 67B \citep{deepseek-llm}, (2) Qwen 72B \citep{qwen}, (3) SeaLLM-v2 7B \citep{seallm}, and (4) ChatGLM3 6B \citep{chatglm3}. Additionally, there are models specifically optimized for mathematical tasks: (5) InternLM2-Math 20B~\footnote{\url{https://github.com/InternLM/InternLM-Math}} extends InternLM2 with specialized math training followed by instruction fine-tuning, (6) Math-Shepherd-Mistral 7B leverages PPO training \citep{schulman2017proximal} on the Mistral 7B \citep{mistral} framework using a process-supervised reward mechanism, (7) the WizardMath family \citep{wizardmath}, which enhances the mathematical performance of both Mistral 7B and Llama-2 70B \citep{llama2} through evolve-instruct (an instruction tuning method with AI-generated instructions) and PPO, using training datasets centered on GSM8K and MATH, (8) MetaMath 70B \citep{metamath}, formed by fine-tuning Llama-2 70B on an expanded dataset consisting of GSM8K and MATH, (9) ToRA 34B \cite{tora}, a version of CodeLlama 34B refined for tool-assisted mathematical reasoning, and (10) MAmmoTH 70B \citep{MathInstruct}, which results from instruction tuning Llama-2 70B with MathInstruct.

As presented in Table \ref{tab:sft_rl_math}, when evaluated without access to external tools, DeepSeekMath-Instruct 7B exhibits robust step-by-step reasoning capabilities. Specifically, on the challenging MATH dataset at the competition level, our model outperforms all open-source alternatives and most closed-source systems (such as Inflection-2 and Gemini Pro) by a margin of no less than 9\% in absolute terms. This superior result holds even when compared to much larger models (for instance, Qwen 72B) or those that have undergone dedicated mathematical reinforcement learning tuning (such as WizardMath-v1.1 7B). Although DeepSeekMath-Instruct performs on par with Chinese proprietary models like GLM-4 and Baichuan-3 in this context, it does not reach the performance level of GPT-4 and Gemini Ultra.

When evaluated under conditions that permit both natural language reasoning and the incorporation of programmatic tools for solving problems, DeepSeekMath-Instruct 7B nearly achieves 60\% accuracy on the MATH dataset, outperforming every currently available open-source alternative. On additional benchmarks, our model demonstrates competitive performance relative to DeepSeek-LLM-Chat 67B, the previous state-of-the-art model that boasts a tenfold increase in scale.

\subsection{Group Relative Policy Optimization} Reinforcement learning (RL) has demonstrated significant success in enhancing the mathematical reasoning skills of LLMs beyond what is achieved through Supervised Fine-Tuning (SFT) \citep{wang2023math,wizardmath}. Here, we present Group Relative Policy Optimization (GRPO), our novel RL algorithm designed for both efficiency and effectiveness.

\subsubsection{From PPO to GRPO} Proximal Policy Optimization (PPO) \citep{schulman2017proximal} is an actor-critic RL algorithm that is widely used in the RL fine-tuning stage of LLMs \citep{ouyang2022training}. In particular, it optimizes LLMs by maximizing the following surrogate objective:

\begin{equation}
\footnotesize
    \mathcal{J}_{PPO}(\theta) = \mathbb{E}{[q \sim P(Q), o \sim \pi_{\theta_{old}}(O|q)]} \frac{1}{|o|} \sum_{t=1}^{|o|} \min \left[ \frac{\pi_\theta(o_{t} | q, o_{<t})}{\pi_{\theta_{old}}(o_{t} | q, o_{<t})} A_{t}, \text{clip} \left( \frac{\pi_\theta(o_{t} | q, o_{<t})}{\pi_{\theta_{old}}(o_{t} | q, o_{<t})}, 1 - \varepsilon, 1 + \varepsilon \right)  A_{t} \right] ,
\end{equation}

The above equation defines the objective function for Proximal Policy Optimization (PPO), denoted by $\mathcal{J}_{PPO}(\theta)$. The expectation is taken with respect to queries $q$ sampled from $P(Q)$ and output sequences $o$ drawn according to the old policy $\pi_{\theta_{old}}(O|q)$. For each timestep $t$ within the output $o$ of length $|o|$, the objective aggregates the clipped surrogate advantage over all tokens. Specifically, at each position, it computes the minimum between the product of the policy ratio $\frac{\pi_\theta(o_{t} | q, o_{<t})}{\pi_{\theta_{old}}(o_{t} | q, o_{<t})}$ and advantage $A_t$, and the advantage $A_t$ scaled by the policy ratio clipped within the interval $[1 - \varepsilon, 1 + \varepsilon]$. Finally, the result is averaged across the sequence length.

In this context, $\pi_{\theta}$ refers to the current policy model, while $\pi_{\theta_{old}}$ denotes the preceding policy version. The variables $q$ and $o$ correspond to the questions and outputs, drawn from both the question dataset and sampled using the old policy $\pi_{\theta_{old}}$. The scalar $\varepsilon$ serves as a clipping hyper-parameter, a feature incorporated in PPO to enhance the reliability of training. The term $A_t$ represents the advantage, which is estimated through the application of Generalized Advantage Estimation (GAE) \citep{gae}, utilizing the sequence of rewards $\{r_{\ge t}\}$ together with an associated value function $V_{\psi}$ that is learned. Consequently, PPO requires joint training of a value function alongside the primary policy model. To avoid excessive optimization directed towards the reward model, it is customary to include, for each token, a KL penalty between the policy and a fixed reference model within the reward structure, as outlined in \citep{ouyang2022training}. Specifically,

\begin{equation}
    r_{t} = r_\varphi(q, o_{\le t}) - \beta \log\frac{\pi_{\theta}(o_{t}|q, o_{<t})}{\pi_{ref}(o_{t}|q, o_{<t})},
\label{eq:PPO-reward}
\end{equation}

Here, $r_\varphi$ denotes the reward model, $\pi_{ref}$ represents the reference model—commonly the original SFT model—and $\beta$ serves as the weighting factor for the KL divergence penalty.

As the value function employed in PPO is typically another model of comparable size as the policy model, it brings a substantial memory and computational burden. Additionally, during RL training, the value function is treated as a baseline in the calculation of the advantage for variance reduction. While in the LLM context, usually only the last token is assigned a reward score by the reward model, which may complicate the training of a value function that is accurate at each token. To address this, as shown in Figure \ref{fig:grpo}, we propose Group Relative Policy Optimization (GRPO), which obviates the need for additional value function approximation as in PPO, and instead uses the average reward of multiple sampled outputs, produced in response to the same question, as the baseline. More specifically, for each question $q$, GRPO samples a group of outputs $\{o_1, o_2, \cdots, o_G\}$  from the old policy  $\pi_{\theta_{old}}$  and then optimizes the policy model by maximizing the following objective:

\begin{equation}
\footnotesize
\begin{split}
    \mathcal{J}_{GRPO}(\theta) &= \mathbb{E}{[q \sim P(Q), \{o_i\}_{i=1}^G \sim \pi_{\theta_{old}}(O|q)]}  \\
    & \frac{1}{G}\sum_{i=1}^G\frac{1}{|o_i|} \sum_{t=1}^{|o_i|} \left\{ \min \left[ \frac{\pi_\theta(o_{i,t} | q, o_{i,<t})}{\pi_{\theta_{old}}(o_{i,t} | q, o_{i,<t})} \hat{A}_{i,t}, \text{clip} \left( \frac{\pi_\theta(o_{i,t} | q, o_{i,<t})}{\pi_{\theta_{old}}(o_{i,t} | q, o_{i,<t})}, 1 - \varepsilon, 1 + \varepsilon \right)  \hat{A}_{i,t} \right] - \beta \mathbb{D}_{KL}\left[\pi_{\theta} || \pi_{ref}\right]\right\} ,
\end{split}
\label{eq:GRPO-obj}
\end{equation}

where $\varepsilon$ and $\beta$ are hyper-parameters, and $\hat{A}_{i,t}$  is the advantage calculated based on relative rewards of the outputs inside each group only, which will be detailed in the following subsections. The group relative way that GRPO leverages to calculate the advantages, aligns well with the comparative nature of rewards models,  as reward models are typically trained on datasets of comparisons between outputs on the same question. Also note that, instead of adding KL penalty in the reward, GRPO regularizes by directly adding the KL divergence between the trained policy and the reference policy to the loss, avoiding complicating the calculation of  $\hat{A}_{i,t}$. And different from the KL penalty term used in (\ref{eq:PPO-reward}), we estimate the KL divergence with the following unbiased estimator \citep{kl_approx}:

\begin{equation}
\small
    \mathbb{D}_{KL}\left[\pi_{\theta} || \pi_{ref}\right] = \frac{\pi_{ref}(o_{i,t}|q,o_{i,<t})}{\pi_{\theta}(o_{i,t}|q,o_{i,<t})} - \log\left(\frac{\pi_{ref}(o_{i,t}|q,o_{i,<t})}{\pi_{\theta}(o_{i,t}|q,o_{i,<t})}\right) - 1,
\end{equation}

which is guaranteed to be positive.

\begin{algorithm}[t]
  \small
  \caption{Iterative Group Relative Policy Optimization}
  \textbf{Input} initial policy model $\pi_{\theta_{\text{init}}}$; reward models $r_\varphi$; task prompts $\mathcal{D}$; 
  hyperparameters $\varepsilon$, $\beta$, $\mu$
  \begin{algorithmic}[1]
    \State Initialize policy model $\pi_\theta$ with $\pi_{\theta_{\text{init}}}$
    \For{iteration = 1 to I}
      \State Set the reference model $\pi_{ref}$ to the current policy $\pi_{\theta}$
      \For{step = 1 to M}
        \State Draw a mini-batch $\mathcal{D}_b$ from $\mathcal{D}$
        \State Assign $\pi_{\theta_{old}}$ as the current policy $\pi_{\theta}$
        \State For every question $q$ in $\mathcal{D}_b$, generate $G$ outputs $\{o_i\}_{i=1}^G$ by sampling from $\pi_{\theta_{old}}(\cdot|q)$
        \State Evaluate each output $o_i$ using the reward model $r_{\varphi}$ to get rewards $\{r_i\}_{i=1}^{G}$
        \State Perform group-relative advantage estimation to obtain $\hat{A}_{i,t}$ for the $t$-th token of each $o_i$
        \For{GRPO iteration = 1 to $\mu$}
          \State Refine the policy $\pi_\theta$ to optimize the GRPO objective as detailed in Equation \ref{eq:GC-GRPO}
        \EndFor
      \EndFor
      \State Continuously train $r_\varphi$ via a replay strategy for updating.
    \EndFor
  \end{algorithmic}
  \textbf{Output} $\pi_\theta$
  \label{alg:iter-grpo}
\end{algorithm}

\subsubsection{Outcome Supervision RL with GRPO} More specifically, for a given question $q$, a collection of responses $\{o_1, o_2, \cdots, o_G\}$ is drawn from the previous iteration of the policy model, denoted as $\pi_{\theta_{old}}$. Each sampled response is evaluated by a reward model, producing a corresponding set of $G$ reward values $\mathbf{r} = \{r_1, r_2, \cdots, r_G\}$. The rewards are then standardized by removing the mean of the group and dividing by the standard deviation. In outcome supervision, the terminal normalized reward for each output $o_i$ is supplied, and this same value is assigned as the advantage $\hat{A}_{i, t}$ to every token within $o_i$; that is, $\hat{A}_{i, t} = \widetilde{r}_i = \frac{r_i- {\rm mean}(\mathbf{r})}{{\rm std}(\mathbf{r})}$. The policy is updated by maximizing the objective outlined in equation (\ref{eq:GRPO-obj}).

\subsubsection{Process Supervision RL with GRPO} Supervising solely via outcome rewards—where feedback is provided only at the end of the generated output—can prove inadequate for guiding the policy during intricate mathematical problem solving. Building on the approach presented by \cite{wang2023math}, we examine process supervision, in which each reasoning step receives an explicit reward signal. More precisely, given a question $q$ and $G$ generated outputs $\{o_1, o_2, \cdots, o_G\}$, a process reward model evaluates each intermediate step, assigning step-specific rewards in the following structure: $\mathbf{R} = \{ \{r_1^{index(1)},\cdots,r_1^{index(K_1)}\}, \cdots,  \{r_G^{index(1)},\cdots,r_G^{index(K_G)}\} \}$. Here, $index(j)$ refers to the token index marking the end of the $j$-th reasoning step, and $K_i$ signifies the step count within the $i$-th generated output. To standardize reward values, we normalize each reward by subtracting the mean of all rewards and dividing by their standard deviation, formalized as $\widetilde{r}_i^{index(j)} = \frac{r_i^{index(j)} - {\rm mean(\mathbf{R})}}{{\rm std(\mathbf{R})}}$.

The process supervision method then computes the advantage for each token as the cumulative sum of these normalized rewards from the current token onward, specifically, $\hat{A}_{i, t} = \sum_{index(j) \ge t} \widetilde{r}_i^{index(j)}$. The policy is updated to maximize the process-level objective as defined in equation (\ref{eq:GRPO-obj}).

\subsubsection{Iterative RL with GRPO} During reinforcement learning, as training advances, the previously trained reward model may no longer be adequate to effectively guide the policy model. To address this, we investigate an iterative approach to RL using GRPO. In this method, outlined in Algorithm \ref{alg:iter-grpo}, updated training datasets for the reward model are obtained by sampling from the current policy model. The reward model is then continually refined, leveraging a replay strategy that maintains 10\% of earlier training data. Afterward, the policy model becomes the new reference model, and further training is conducted on the policy model, this time employing the updated reward model for guidance.

\subsection{Training and Evaluating DeepSeekMath-RL}

Our reinforcement learning experiments utilize DeepSeekMath-Instruct 7B as the foundation. The RL training data comprises chain-of-thought-style problems drawn from GSM8K and MATH sections within the SFT dataset, totaling approximately 144,000 questions. We purposefully omit all other SFT questions to isolate and analyze the influence of RL on evaluation sets lacking overlap with the RL data. For constructing the reward model training set, we replicate the approach described in \citep{wang2023math}. The initial reward model is trained on DeepSeekMath-Base 7B using a learning rate of $2 \times 10^{-5}$. In the GRPO framework, we assign the policy model a learning rate of $1 \times 10^{-6}$ and fix the KL regularization coefficient at 0.04. For every question, we generate $64$ candidate responses. The model’s maximum sequence length and batch size are both set to 1024. After each exploratory episode, the policy model undergoes a single parameter update. We assess the performance of DeepSeekMath-RL 7B using the same benchmark protocols established for DeepSeekMath-Instruct 7B. Within this evaluation scheme, chain-of-thought tasks in GSM8K and MATH are treated as in-domain for DeepSeekMath-RL 7B, while the remaining benchmarks serve as out-of-domain assessments.

Table \ref{tab:sft_rl_math} presents comparative results for both open-source and proprietary models that employ chain-of-thought and tool-augmented reasoning strategies on English and Chinese evaluation tasks. The key observations are as follows: 1) DeepSeekMath-RL 7B achieves accuracy rates of 88.2\% on GSM8K and 51.7\% on MATH with chain-of-thought reasoning, exceeding the performance of all open-source counterparts within the 7B to 70B scale and outperforming most closed-source alternatives. 2) Importantly, the training regimen for DeepSeekMath-RL 7B relies exclusively on chain-of-thought style instruction-tuning datasets drawn from GSM8K and MATH, using DeepSeekMath-Instruct 7B as its starting point. Although its training data is narrowly focused, DeepSeekMath-RL 7B achieves higher scores than DeepSeekMath-Instruct 7B across all tested criteria, thereby demonstrating the advantages conferred by reinforcement learning.

\section{Discussion}
This section presents an overview of the results obtained from our pre-training and reinforcement learning experiments.

\subsection{Lessons Learnt in Pre-Training}

We begin by presenting insights gained during the pre-training stage. Unless explicitly mentioned, our training procedures follow those described in Section \ref{sec:quality-policy}. Additionally, in this context, all references to the DeepSeekMath Corpus pertain to the 89B-token dataset assembled during the second round of data collection.

\subsubsection{Code Training Benefits Mathematical Reasoning}

A widely held but largely untested assumption posits that learning to code enhances reasoning skills. In this work, we address this claim in the context of mathematics, providing evidence that code training strengthens models’ mathematical reasoning capabilities, irrespective of tool usage.

To study how code training affects mathematical reasoning, we experimented with the following two-stage training and one-stage training settings:

\textbf{Two-Stage Training}

\begin{itemize}[topsep=0pt]
    \item \textbf{Code Training for 400B Tokens $\rightarrow$ Math Training for 150B Tokens}: DeepSeek-LLM 1.3B is initially pretrained on 400B tokens of code data, after which it undergoes an additional 150B tokens of mathematics-focused training;
    \item \textbf{General Training for 400B Tokens $\rightarrow$ Math Training for 150B Tokens}: For comparison, we run a baseline where the model is exposed to 400B general-domain tokens (sourced from DeepSeek-AI’s extensive general corpus) during the first phase, aiming to analyze the comparative benefits of code-centric pretraining as opposed to general-domain pretraining for mathematical reasoning capability.
\end{itemize}

\textbf{One-Stage Training}

\begin{itemize}[topsep=0pt]
    \item \textbf{Math Pretraining on 150B Tokens}: DeepSeek-LLM 1.3B is pretrained using a dataset comprising 150 billion math tokens;
    \item \textbf{Joint Training with 400B Code Tokens and 150B Math Tokens}: Sequential math training after code training leads to diminished code proficiency. We examine if integrating code and math tokens in a single unified training phase can simultaneously bolster mathematical reasoning capabilities and mitigate catastrophic forgetting.
\end{itemize}

\paragraph{Results}
The outcomes presented in Table \ref{tab:code-math} and Table \ref{tab:code-math-general-eval} illustrate the downstream effectiveness across various training configurations.

Incorporating code training supports program-aided mathematical reasoning, both in two-stage and one-stage training paradigms. As evidenced in Table \ref{tab:code-math}, within the two-stage framework, exposure to code training alone substantially boosts the model’s capacity to address GSM8K and MATH tasks via Python. Introducing mathematical reasoning tasks in the subsequent training stage leads to additional gains. Notably, in the one-stage training setup, integrating code and math tokens within the same process not only prevents the catastrophic forgetting often observed in two-stage training, but also enhances performance in coding (see Table \ref{tab:code-math-general-eval}) and in tasks requiring program-facilitated mathematical reasoning (refer to Table \ref{tab:code-math}).

Code training independently enhances mathematical reasoning skills even in the absence of tool use. In the two-stage training approach, moderate gains are observable following the preliminary phase dedicated to code training. Furthermore, this phase facilitates more effective subsequent mathematical training, culminating in optimal performance. In contrast, merging code and math tokens into a unified one-stage training regimen undermines mathematical reasoning when tool use is excluded. A possible explanation is that DeepSeek-LLM 1.3B may not possess sufficient scale to adequately encode both code-related and mathematical information concurrently.

\subsubsection{ArXiv Papers Seem Ineffective in Improving Mathematical Reasoning}

ArXiv papers are commonly included as a component of math pre-training data \citep{minerva,gpt-f,llemma,mathpile}. However, detailed analysis regarding their impact on mathematical reasoning has not been extensively conducted. Perhaps counter-intuitively, according to our experiments, arXiv papers seem ineffective in improving mathematical reasoning. We experiment with models of different sizes, including DeepSeek-LLM 1.3B and DeepSeek-Coder-Base-v1.5 7B \citep{deepseek-coder}, using arXiv corpora that underwent varied processing pipelines:

\begin{itemize}[topsep=0pt]
    \item \textbf{MathPile} \citep{mathpile}: a dataset composed of 8.9 billion tokens, assembled using various heuristic filtering and cleaning methods, where more than 85\% of its content originates from scientific papers on arXiv;
    \item \textbf{ArXiv-RedPajama} \citep{redpajama}: a dataset consisting of all arXiv LaTeX documents, excluding preambles, comments, macros, and reference sections, amounting to 28.0 billion tokens in total.
\end{itemize}

For our study, we independently pre-trained DeepSeek-LLM 1.3B using 150B tokens and DeepSeek-Coder-Base-v1.5 7B with 40B tokens from each arXiv dataset. Our findings indicate that incorporating arXiv papers does not meaningfully enhance models’ mathematical reasoning abilities. Training solely on arXiv data, neither model achieved significant gains and, in some cases, even demonstrated reduced performance across a range of mathematical benchmarks that varied in complexity. These benchmarks included datasets for quantitative reasoning such as GSM8K and MATH (Table \ref{tab:arxiv-cot}), multiple-choice evaluations like MMLU-STEM (Table \ref{tab:arxiv-cot}), and formal mathematics tasks, represented by miniF2F (Table \ref{tab:arxiv_atp}).

However, this conclusion has its limitations and should be taken with a grain of salt. We have not yet studied:

\begin{itemize}[topsep=0pt]
    \item How arXiv tokens influence math-centered tasks beyond those examined in this study, for instance, transforming formal mathematical statements or proofs into informal equivalents;
    \item The influence of arXiv tokens in conjunction with alternative data modalities;
    \item If scaling up models would reveal advantages associated with incorporating arXiv papers.
\end{itemize}

Thus, further exploration is required, which we leave for future studies.

\subsection{Insights of Reinforcement Learning}
\subsubsection{Towards to a Unified Paradigm} Here, we introduce a comprehensive framework to systematically examine a variety of training approaches, including SFT, RFT, DPO, PPO, and GRPO. Moreover, we empirically investigate components influencing this unified framework. In general terms, the gradient of a training method with respect to its parameter $\theta$ can be formulated as follows:

\begin{equation}
    \nabla_{\theta}\mathcal{J}_{\textcolor{red}{\mathcal{A}}}(\theta) = \mathbb{E}[\underbrace{(q,o) \sim \textcolor{red}{\mathcal{D}}}_{Data \ Source}]\left( \frac{1}{|o|} \sum_{t=1}^{|o|}  \underbrace{GC_{{\mathcal{A}}}(q, o, t, \textcolor{red}{\pi_{{rf}}})}_{Gradient \ Coefficient}  \nabla_{\theta}\log \pi_{\theta}(o_t | q, o_{<t})\right).
\label{eq:objective}
\end{equation}

There exist three key components: 1) \textit{Data Source $\mathcal{D}$}, which determines the training data; 2) \textit{Reward Function $\pi_{{rf}}$}, which is the source of the training reward signal; 3) \textit{Algorithm $\mathcal{A}$}: which processes the training data and the reward signal to the gradient coefficient $GC$ that determines the magnitude of the penalty or reinforcement for the data. We analyze several representative methods based on such a unified paradigm:

\begin{itemize}
    \item \textbf{Supervised Fine-tuning (SFT)}: SFT involves adapting a pretrained model by training it on a curated dataset selected by humans.
    \item \textbf{Rejection Sampling Fine-tuning (RFT)}: RFT enhances the SFT model by additional fine-tuning on outputs generated by the SFT model in response to SFT prompts; these outputs are filtered to include only those with correct answers.
    \item \textbf{Direct Preference Optimization (DPO)}: DPO builds upon the SFT model, further training it using an augmented set of outputs from the SFT model and applying a pair-wise DPO loss function to optimize preferences.
    \item \textbf{Online Rejection Sampling Fine-tuning (Online RFT)}: Unlike RFT, Online RFT initializes its policy from the SFT model and continuously updates it through fine-tuning on outputs that are generated and filtered from the current policy in real-time.
    \item \textbf{PPO/GRPO}: PPO/GRPO methods start from the SFT model as the policy initialization and apply reinforcement by training on outputs sampled on-the-fly from the up-to-date policy model.
\end{itemize}

Table \ref{tab:data_policy} provides a summary of the components for these methods. For a comprehensive explanation of the derivation, see Appendix \ref{app:analysis-rl}.

\paragraph{Observation about Data Source} The data source can be categorized into two types: online sampling and offline sampling. Online sampling refers to obtaining training data from the outputs generated by the current policy model during live exploration. In contrast, offline sampling refers to collecting training data from the samples generated by the pre-trained SFT model. Both RFT and DPO utilize offline sampling, whereas Online RFT and GRPO adopt the online sampling approach.

As illustrated in Figure \ref{fig:rl-analysis}, our results indicate that Online RFT achieves considerably better performance than RFT across two benchmark tasks. In particular, while Online RFT and RFT perform similarly during the initial training phase, Online RFT displays clear superiority in the subsequent stages, underscoring the benefits of online learning. This observation is reasonable: during the early stages, there is little distinction between the outputs of the actor and the SFT model, leading to sampled data that are largely alike. As training advances, the divergence between the actor and the SFT model grows, resulting in increasingly distinct data from the actor, where the advantage of real-time data generation becomes more pronounced.

\paragraph{Observation about Gradient Coefficient} The reward signal is utilized by the algorithm to compute the gradient coefficient for updating model parameters. In our study, we categorize the reward function into two types: ‘Rule’ and ‘Model’. The ‘Rule’ type evaluates the response quality solely by assessing answer correctness, whereas ‘Model’ involves employing a reward model that assigns a score to each response. The reward model’s training data is derived from judgements made using the ‘Rule’ criterion. As shown in Equations \ref{eq:GC-RFT} and \ref{eq:GC-GRPO}, there exists a fundamental distinction between GRPO and Online RFT: GRPO modifies its gradient coefficient dynamically in accordance with the reward model’s output, facilitating nuanced reinforcement and penalty depending on reward magnitude. By comparison, Online RFT does not possess this capability; it abstains from penalizing wrong responses and applies equal reinforcement to all responses deemed correct, regardless of their reward values.

As illustrated in Figure \ref{fig:rl-analysis}, GRPO outperforms online RFT, underscoring the effectiveness of adjusting the coefficients for positive and negative gradients. Moreover, GRPO+PS achieves better results than GRPO+OS, which suggests that employing more granular, step-specific gradient coefficients is advantageous. Additionally, we investigate iterative RL by carrying out two iterations in our experiments. Figure \ref{fig:iter-rl} demonstrates that applying iterative RL leads to marked improvements in performance, particularly following the initial iteration.

\subsubsection{Why RL Works?} In this work, we apply reinforcement learning utilizing a selected portion of the instruction tuning dataset, observing notable gains in performance over the standard instruction tuning approach. To investigate the underlying mechanisms responsible for these improvements, we compare the Pass@K and Maj@K accuracy metrics for both the Instruct and RL-based models across two evaluation benchmarks. As depicted in Figure \ref{fig:maj-pass}, reinforcement learning enhances Maj@K scores while leaving Pass@K largely unaffected. This pattern suggests that the effectiveness of RL mainly lies in improving the reliability of the model's output distribution—specifically, \textbf{the observed gains likely result from an increased likelihood of returning the correct answer among the TopK candidates, rather than advancing the model’s core abilities.} Analogously, \citep{wang2023making} reported a \textbf{misalignment problem} when dealing with reasoning tasks in SFT models, and further demonstrated that strategies involving preference alignment can systematically improve SFT model reasoning outcomes \citep{yuan2023rrhf,song2023preference,wang2023making}.

\subsubsection{How to Achieve More Effective RL?}
\label{sec:effec_rl}
Our results indicate that RL is highly effective for mathematical reasoning tasks. Furthermore, we present a cohesive framework that elucidates the commonalities among various notable training approaches, recasting them as either explicit or approximate RL strategies. This unified view, encapsulated in Equation \ref{eq:objective}, distinguishes three principal elements: Data Source, Algorithm, and Reward Function. We discuss possible avenues for advancing these components in future work.

\paragraph{Data Source} The foundation of all training procedures lies in the data source. For RL, the data source is defined here as the unlabeled prompts alongside their corresponding outputs, which are generated by sampling from the policy model. In our experiments, we limit ourselves to prompts utilized during the instruction tuning phase and employ a basic nucleus sampling scheme to generate responses. We believe this limitation could contribute to why our RL pipeline yields improvements solely in Maj@K metrics. Looking ahead, we intend to assess our RL approach on question prompts that fall outside the training distribution, as well as experiment with \textbf{advanced sampling (decoding) strategies} informed by techniques such as tree-search methods \citep{yao2023tree}. Moreover, innovations in \textbf{efficient inference techniques} \citep{xia-etal-2023-speculative,leviathan2023fast,kwon2023efficient,xia2024unlocking}—which shape how effectively the policy model explores—are expected to significantly impact performance.

\paragraph{Algorithms} Algorithms utilize both the observed data and the reward signal to compute gradient coefficients, thereby facilitating the update of model parameters. According to Equation \ref{eq:objective}, current methodologies generally place complete \textbf{TRUST} in the reward function's feedback when adjusting the conditional probability of specific tokens. Nevertheless, the reliability of the reward signal cannot be guaranteed in all scenarios, particularly when dealing with highly intricate tasks. For instance, even datasets such as PRM800K \citep{lightman2023let}, which have undergone thorough annotation by expert annotators, still exhibit an error rate of around 20\% in their labels\footnote{\url{https://github.com/openai/prm800k/issues/12\#issuecomment-1728491852}}. Consequently, we intend to investigate reinforcement learning algorithms that maintain robustness in the presence of noisy reward feedback. Our hypothesis is that alignment strategies following the \textbf{WEAK-TO-STRONG} \citep{burns2023weak} paradigm will fundamentally transform the landscape of learning algorithms.

\paragraph{Reward Function} The reward function provides the essential training feedback. In RL settings, this function is commonly represented by a neural reward model. We identify three key areas of research regarding reward models: 1) \textbf{Improving the generalization capacity of reward models.} It is crucial for reward models to generalize adequately to previously unseen questions and novel outputs generated during decoding; failure to do so risks reinforcement learning methods serving only to reinforce the current distribution of LLMs without yielding substantive improvement in their core skills; 2) \textbf{Incorporating the reward model’s uncertainty.} Quantifying uncertainty may serve as a bridge between limited reward model capabilities and algorithms aimed at bootstrapping weaker models into stronger ones; 3) \textbf{Developing efficient methods for constructing high-quality process reward models} capable of delivering granular training feedback throughout the reasoning steps \citep{lightman2023let,wang2023math}.

\section{Conclusion, Limitation, and Future Work} In this work, we introduce DeepSeekMath, a model that surpasses all open-source counterparts on the competitive MATH benchmark and nearly matches the performance observed in proprietary models. DeepSeekMath is based on DeepSeek-Coder-v1.5 7B as its starting point and is subjected to continued pretraining for 500B tokens, where a substantial portion, specifically 120B math-specific tokens, is derived from Common Crawl. Our detailed ablation analyses highlight that web content serves as a valuable reservoir for high-quality mathematical data, whereas arXiv does not confer as much advantage as initially anticipated. We propose Group Relative Policy Optimization (GRPO), an adaptation of Proximal Policy Optimization (PPO), which offers marked enhancements to mathematical reasoning proficiency while reducing memory requirements. Empirical findings indicate the robustness of GRPO, with effectiveness persisting even when DeepSeekMath-Instruct 7B achieves elevated scores across benchmarks. Additionally, we present a unified framework for conceptualizing related methods and outline several promising research avenues to further improve reinforcement learning strategies.

While DeepSeekMath~demonstrates strong results on various quantitative reasoning tasks, it lags behind proprietary models in areas such as geometry and theorem proving. For example, in preliminary testing, the model struggled with questions involving triangles and ellipses, which could reflect bias in the data used for both pre-training and fine-tuning stages. Furthermore, due to limitations in model size, DeepSeekMath~falls short of GPT-4 in terms of few-shot learning; whereas GPT-4 displays clear improvement with few-shot examples, DeepSeekMath~performs comparably on both zero-shot and few-shot tasks. Moving forward, we plan to enhance our engineered data selection process to develop a higher-quality pre-training corpus. Additionally, we intend to investigate alternative approaches (see Section~\ref{sec:effec_rl}) for enabling more effective reinforcement learning with LLMs.

\end{document}